\section{Metodología de Evaluación}

Para evaluar con precisión los resultados y el impacto de la implementación de Odoo en PERÚ FLORES TRAVEL AGENCY TOURS OPERADOR E.I.R.L., se aplicó un enfoque de análisis integral orientado a medir la transición desde procesos manuales hacia un ecosistema digital centralizado. La evaluación se estructuró bajo los siguientes cuatro pilares:

\begin{itemize}
    \item \textbf{Comparación de procesos (Antes vs. Después):} Se realizó un análisis comparativo entre el modelo operativo previo ---basado en el uso fragmentado de Excel, WhatsApp y cuadernos físicos--- y el nuevo entorno digital integrado en Odoo. Esta evaluación permitió identificar la eliminación de "silos de información" y la mejora en la trazabilidad de los pasajeros, desde el primer contacto comercial hasta la ejecución del servicio.
    \item \textbf{Análisis de eficiencia operativa:} Se evaluó el desempeño de los procesos directamente en el sistema, verificando la automatización de las cotizaciones de paquetes turísticos (Cusco, Arequipa, Puno). El enfoque se centró en medir la reducción de tiempos de respuesta al cliente, la estandarización de itinerarios y la disminución de errores manuales en la generación de presupuestos en PDF.
    \item \textbf{Evaluación de Indicadores de Desempeño (KPIs):} Se definieron indicadores clave basados en el diagnóstico inicial de 2025 para cuantificar las mejoras. Los indicadores principales analizados fueron el Tiempo Promedio de Cotización (con una meta de reducción a menos de 15 minutos) y la Tasa de Error Operativo en la reserva de servicios (con una meta de reducción por debajo del 1\%).
    \item \textbf{Observación directa y nivel de adopción:} Se realizó un seguimiento del uso del sistema en el entorno real de trabajo, evaluando la interacción de la gerencia y el personal con la plataforma tanto en escritorio como en la aplicación móvil. Este análisis permitió medir la reducción en el uso de herramientas externas no integradas y la correcta ejecución de los flujos de trabajo dentro del ERP.
\end{itemize}
